% Options for packages loaded elsewhere
% Options for packages loaded elsewhere
\PassOptionsToPackage{unicode}{hyperref}
\PassOptionsToPackage{hyphens}{url}
\PassOptionsToPackage{dvipsnames,svgnames,x11names}{xcolor}
%
\documentclass[
  japanese,
  lualatex,
  ja=standard,
  12pt,
  a4paper]{ltjsarticle}
\usepackage{xcolor}
\usepackage{amsmath,amssymb}
\setcounter{secnumdepth}{5}
\usepackage{iftex}
\ifPDFTeX
  \usepackage[T1]{fontenc}
  \usepackage[utf8]{inputenc}
  \usepackage{textcomp} % provide euro and other symbols
\else % if luatex or xetex
  \usepackage{unicode-math} % this also loads fontspec
  \defaultfontfeatures{Scale=MatchLowercase}
  \defaultfontfeatures[\rmfamily]{Ligatures=TeX,Scale=1}
\fi
\usepackage{lmodern}
\ifPDFTeX\else
  % xetex/luatex font selection
\fi
% Use upquote if available, for straight quotes in verbatim environments
\IfFileExists{upquote.sty}{\usepackage{upquote}}{}
\IfFileExists{microtype.sty}{% use microtype if available
  \usepackage[]{microtype}
  \UseMicrotypeSet[protrusion]{basicmath} % disable protrusion for tt fonts
}{}
\makeatletter
\@ifundefined{KOMAClassName}{% if non-KOMA class
  \IfFileExists{parskip.sty}{%
    \usepackage{parskip}
  }{% else
    \setlength{\parindent}{0pt}
    \setlength{\parskip}{6pt plus 2pt minus 1pt}}
}{% if KOMA class
  \KOMAoptions{parskip=half}}
\makeatother
% Make \paragraph and \subparagraph free-standing
\makeatletter
\ifx\paragraph\undefined\else
  \let\oldparagraph\paragraph
  \renewcommand{\paragraph}{
    \@ifstar
      \xxxParagraphStar
      \xxxParagraphNoStar
  }
  \newcommand{\xxxParagraphStar}[1]{\oldparagraph*{#1}\mbox{}}
  \newcommand{\xxxParagraphNoStar}[1]{\oldparagraph{#1}\mbox{}}
\fi
\ifx\subparagraph\undefined\else
  \let\oldsubparagraph\subparagraph
  \renewcommand{\subparagraph}{
    \@ifstar
      \xxxSubParagraphStar
      \xxxSubParagraphNoStar
  }
  \newcommand{\xxxSubParagraphStar}[1]{\oldsubparagraph*{#1}\mbox{}}
  \newcommand{\xxxSubParagraphNoStar}[1]{\oldsubparagraph{#1}\mbox{}}
\fi
\makeatother


\usepackage{longtable,booktabs,array}
\newcounter{none} % for unnumbered tables
\usepackage{calc} % for calculating minipage widths
% Correct order of tables after \paragraph or \subparagraph
\usepackage{etoolbox}
\makeatletter
\patchcmd\longtable{\par}{\if@noskipsec\mbox{}\fi\par}{}{}
\makeatother
% Allow footnotes in longtable head/foot
\IfFileExists{footnotehyper.sty}{\usepackage{footnotehyper}}{\usepackage{footnote}}
\makesavenoteenv{longtable}
\usepackage{graphicx}
\makeatletter
\newsavebox\pandoc@box
\newcommand*\pandocbounded[1]{% scales image to fit in text height/width
  \sbox\pandoc@box{#1}%
  \Gscale@div\@tempa{\textheight}{\dimexpr\ht\pandoc@box+\dp\pandoc@box\relax}%
  \Gscale@div\@tempb{\linewidth}{\wd\pandoc@box}%
  \ifdim\@tempb\p@<\@tempa\p@\let\@tempa\@tempb\fi% select the smaller of both
  \ifdim\@tempa\p@<\p@\scalebox{\@tempa}{\usebox\pandoc@box}%
  \else\usebox{\pandoc@box}%
  \fi%
}
% Set default figure placement to htbp
\def\fps@figure{htbp}
\makeatother



\ifLuaTeX
\usepackage[bidi=basic,shorthands=off,provide=*]{babel}
\else
\usepackage[bidi=default,shorthands=off,provide=*]{babel}
\fi
\ifLuaTeX
  \usepackage{selnolig} % disable illegal ligatures
\fi


\setlength{\emergencystretch}{3em} % prevent overfull lines

\providecommand{\tightlist}{%
  \setlength{\itemsep}{0pt}\setlength{\parskip}{0pt}}



 


\usepackage{needspace}
\makeatletter
\@ifpackageloaded{caption}{}{\usepackage{caption}}
\AtBeginDocument{%
\ifdefined\contentsname
  \renewcommand*\contentsname{目次}
\else
  \newcommand\contentsname{目次}
\fi
\ifdefined\listfigurename
  \renewcommand*\listfigurename{図一覧}
\else
  \newcommand\listfigurename{図一覧}
\fi
\ifdefined\listtablename
  \renewcommand*\listtablename{表一覧}
\else
  \newcommand\listtablename{表一覧}
\fi
\ifdefined\figurename
  \renewcommand*\figurename{図}
\else
  \newcommand\figurename{図}
\fi
\ifdefined\tablename
  \renewcommand*\tablename{表}
\else
  \newcommand\tablename{表}
\fi
}
\@ifpackageloaded{float}{}{\usepackage{float}}
\floatstyle{ruled}
\@ifundefined{c@chapter}{\newfloat{codelisting}{h}{lop}}{\newfloat{codelisting}{h}{lop}[chapter]}
\floatname{codelisting}{コード}
\newcommand*\listoflistings{\listof{codelisting}{コード一覧}}
\makeatother
\makeatletter
\makeatother
\makeatletter
\@ifpackageloaded{caption}{}{\usepackage{caption}}
\@ifpackageloaded{subcaption}{}{\usepackage{subcaption}}
\makeatother
\usepackage{bookmark}
\IfFileExists{xurl.sty}{\usepackage{xurl}}{} % add URL line breaks if available
\urlstyle{same}
% fallback for those not using the hyperref driver hyperxmp:
\makeatletter
\@ifundefined{xmpquote}{\newcommand{\xmpquote}[1]{#1}}{}
\makeatother
\hypersetup{
  pdftitle={第10回：価格・賃金硬直性とRANKモデル},
  pdfauthor={宮﨑憲治},
  pdflang={ja},
  colorlinks=true,
  linkcolor={blue},
  filecolor={Maroon},
  citecolor={Blue},
  urlcolor={Blue},
  pdfcreator={LaTeX via pandoc}}


\title{第10回：価格・賃金硬直性とRANKモデル}
\usepackage{etoolbox}
\makeatletter
\providecommand{\subtitle}[1]{% add subtitle to \maketitle
  \apptocmd{\@title}{\par {\large #1 \par}}{}{}
}
\makeatother
\subtitle{マクロ経済学B}
\author{宮﨑憲治}
\date{2026年度}
\begin{document}
\maketitle

\renewcommand*\contentsname{目次}
{
\setcounter{tocdepth}{3}
\tableofcontents
}

\section{講義の目的}\label{ux8b1bux7fa9ux306eux76eeux7684}

本講義では、第7回で構築した価格硬直性のみのモデルを基礎とし、そこに\textbf{「名目賃金硬直性（Sticky
Nominal Wages）」}を追加した、\textbf{二重の名目硬直性（Double Nominal
Rigidities）をもつ RANK
モデル}を学びます。なお、家計部門は依然として代表的家計の枠組みを維持しているため、ここで扱うのは後続の講義で学ぶ
TANK や HANK ではなく、代表的家計モデルの自然な拡張としての RANK
モデルです。

本回でも、計算の見通しを最も明快にするため、第7回と同様に資本分配率を \[
\alpha=0
\]
に固定します。すなわち、集計生産関数は労働に関して線形であり、政府支出も導入しません。したがって、本講義で新たに学ぶ核心的な論点は、資本の収穫逓減や財政政策ではなく、「\textbf{財の価格（物価）}と\textbf{労働の価格（名目賃金）}という2つの名目変数がともにゆっくりとしか調整されないとき、ニューケインジアン・モデルの動学体系や実質賃金の調整メカニズムがどのように変容するか」という点にあります。

先に結論を整理すると、価格硬直性のみのモデルと比べた本モデルの主要な相違点は以下の通りです。

{\def\LTcaptype{none} % do not increment counter
\begin{longtable}[]{@{}
  >{\raggedright\arraybackslash}p{(\linewidth - 4\tabcolsep) * \real{0.3333}}
  >{\raggedright\arraybackslash}p{(\linewidth - 4\tabcolsep) * \real{0.3333}}
  >{\raggedright\arraybackslash}p{(\linewidth - 4\tabcolsep) * \real{0.3333}}@{}}
\toprule\noalign{}
\begin{minipage}[b]{\linewidth}\raggedright
比較論点
\end{minipage} & \begin{minipage}[b]{\linewidth}\raggedright
価格硬直性のみの RANK モデル（第7回）
\end{minipage} & \begin{minipage}[b]{\linewidth}\raggedright
価格・賃金硬直性をもつ RANK モデル（第10回）
\end{minipage} \\
\midrule\noalign{}
\endhead
\bottomrule\noalign{}
\endlastfoot
家計の構成 & 代表的家計 &
代表的家計（\textbf{家計内で完全な保険を共有}） \\
労働市場の構造 & 同質な労働、完全に柔軟な実質賃金 &
差別化された労働、労働集計業者（\textbf{労働パッカー}）、硬直的な名目賃金 \\
賃金決定条件 & 賃金マークアップは常にゼロ（\(\mu^w_t=0\)） &
賃金フィリップス曲線（\(\pi^w_t=\beta\mathbb{E}_t\pi^w_{t+1}-\kappa_w\mu^w_t\)） \\
価格インフレ率 & 価格マークアップ（\textbf{実質限界費用}）の乖離に反応 &
価格マークアップの乖離に反応 \\
賃金インフレ率 &
名目賃金は柔軟に変化し、独立した賃金フィリップス曲線は不要 &
賃金マークアップの乖離に反応 \\
実質賃金の調整 & 限界代替率と一致して瞬時に調整 &
賃金インフレ率と価格インフレ率の差を通じて緩慢に調整 \\
社会厚生損失の対象 & 需給ギャップ、価格インフレ率 &
需給ギャップ、価格インフレ率、および賃金インフレ率 \\
\end{longtable}
}

本回では、定常状態における財市場の独占マークアップおよび労働市場の独占マークアップの歪みを、それぞれ売上補助金
\(\tau_p^*\) と賃金補助金 \(\tau_w^*\)
によって完全に相殺します。これにより、価格も賃金も完全に柔軟な場合の自然配分はパレート効率的となります。したがって、
\[
y_t^e=y_t^f,
\qquad
r_t^e=r_t^f,
\qquad
x_t=y_t-y_t^f=y_t-y_t^e
\] が成り立ち、需給ギャップと厚生上の産出ギャップは完全に一致します。

本講義を貫く最も中心的な2本の方程式は、以下の通りです。 \[
\mu_t^p+\mu_t^w=-(\gamma+\varphi)x_t,
\qquad
\pi_t^w-\pi_t^p=\Delta w_t \equiv w_t - w_{t-1}.
\]
第1式は\textbf{「2つのマークアップの和が需給ギャップに直結すること」}を示し、第2式は\textbf{「名目賃金インフレ率と価格インフレ率の差が実質賃金の動学的変化を生み出すこと」}を表しています。

図~\ref{fig-lecture10-overview}
は、本講義の全体的な構造を示しています。価格硬直性は企業の動学的価格設定から生じ、賃金硬直性は家計の動学的賃金設定から生じます。この2つの名目硬直性は、実質賃金の動学を通じて密接に相互連関しています。

\begin{figure}

\centering{

\includegraphics[width=0.95\linewidth,height=\textheight,keepaspectratio]{lecture10_files/mediabag/figures/lecture10/lecture10_overview_infographic.pdf}

}

\caption{\label{fig-lecture10-overview}第10回講義の全体像：第7回の価格硬直性
RANK モデルに労働パッカーと Rotemberg
型賃金調整費用を加えると、\(\mu_t^p+\mu_t^w=-(\gamma+\varphi)x_t\) と
\(\pi_t^w-\pi_t^p=\Delta w_t\)
という2本の要の方程式が価格・賃金ブロックを相互に結合する。}

\end{figure}%

\section{表記と基本設定}\label{ux8868ux8a18ux3068ux57faux672cux8a2dux5b9a}

本ノートにおける表記慣例として、大文字は各変数の「水準（\textbf{Level}）」を表し、小文字は原則としてゼロインフレ定常状態からの「対数乖離（\textbf{Log-deviation}）」を表します。ただし、定常値がゼロの純配当
\(d_t\) は、定常消費で規格化した一次偏差です。第10回では
\(\alpha=0\)、政府支出なし、\(a=0\)、\(\nu=1\)
と設定し、定常状態の技術水準 \(\exp(a)=1\) のもとで \[
N=Y=C=W=1,
\qquad
MU^p=MU^w=1
\] となるように正規化します。

本講義では価格硬直性と賃金硬直性を同時に扱うため、この回に限り、価格に関する変数には上付き記号
\(p\)、賃金に関する変数には上付き記号 \(w\)
を付与して区別します。例えば、他の講義回で \(\pi_t\)
と表記していた価格インフレ率は本回では \(\pi_t^p\)
と表記し、賃金インフレ率を \(\pi_t^w\)
と表記します。同様に、代替弾力性や調整費用のパラメータも
\(\psi_p, \eta_p\) と \(\psi_w, \eta_w\) のように下付き記号 \(p, w\)
で明確に区別します。

主な水準・非線形変数の定義は以下の通りです。

{\def\LTcaptype{none} % do not increment counter
\begin{longtable}[]{@{}ll@{}}
\toprule\noalign{}
記号 & 経済学的な意味 \\
\midrule\noalign{}
\endhead
\bottomrule\noalign{}
\endlastfoot
\(C_t\) & 集計民間消費 \\
\(N_t\) & 集計労働投入量 \\
\(N_t(h)\) & 労働タイプ \(h\) の個別労働供給量 \\
\(Y_t\) & 集計産出量 \\
\(W_t^N\) & 集計名目賃金指数 \\
\(W_t^N(h)\) & 労働タイプ \(h\) が設定する個別名目賃金 \\
\(s_t^w(h)=W_t^N(h)/W_t^N\) & 労働タイプ \(h\) の相対賃金 \\
\(W_t=W_t^N/P_t\) & 集計実質賃金 \\
\(P_t\) & 最終財の物価水準 \\
\(\Pi_t^p=P_t/P_{t-1}\) & 価格の粗インフレ率 \\
\(\Pi_t^w=W_t^N/W_{t-1}^N\) & 賃金の粗インフレ率 \\
\(R_t^N\) & 粗名目利子率 \\
\(Q_{t,t+1}\) & 家計の確率的割引因子（\textbf{SDF}） \\
\(MU_t^p\) & 財市場における価格マークアップ \\
\(MU_t^w\) & 労働市場における賃金マークアップ \\
\(a_t\) & 全要素生産性ショックの対数乖離（\textbf{技術水準は
\(\exp(a_t)\)}） \\
\end{longtable}
}

主な対数乖離変数の定義は以下の通りです。

{\def\LTcaptype{none} % do not increment counter
\begin{longtable}[]{@{}
  >{\raggedright\arraybackslash}p{(\linewidth - 2\tabcolsep) * \real{0.5000}}
  >{\raggedright\arraybackslash}p{(\linewidth - 2\tabcolsep) * \real{0.5000}}@{}}
\toprule\noalign{}
\begin{minipage}[b]{\linewidth}\raggedright
記号
\end{minipage} & \begin{minipage}[b]{\linewidth}\raggedright
経済学的な意味
\end{minipage} \\
\midrule\noalign{}
\endhead
\bottomrule\noalign{}
\endlastfoot
\(c_t, y_t, n_t, w_t\) &
消費、産出、労働投入、実質賃金の定常状態からの対数乖離 \\
\(\pi_t^p=\log\Pi_t^p\) & 価格インフレ率の対数乖離 \\
\(\pi_t^w=\log\Pi_t^w\) & 賃金インフレ率の対数乖離 \\
\(r_t^N=\log(R_t^N/R^N)\) & 粗名目利子率の対数乖離 \\
\(r_t=r_t^N-\mathbb{E}_t\pi^p_{t+1}\) & 事前的実質利子率 \\
\(\mu_t^p\) & 価格マークアップの定常状態からの対数乖離 \\
\(\mu_t^w\) & 賃金マークアップの定常状態からの対数乖離 \\
\(y_t^f=y_t^e\) &
柔軟価格・柔軟賃金下の自然産出量（\textbf{本モデルでは効率的産出量に一致}） \\
\(r_t^f=r_t^e\) &
自然実質利子率（\textbf{本モデルでは効率的利子率に一致}） \\
\(x_t=y_t-y_t^f=y_t-y_t^e\) &
需給ギャップ（\textbf{厚生上の産出ギャップ}） \\
\end{longtable}
}

主なパラメータの設定は以下の通りです。

{\def\LTcaptype{none} % do not increment counter
\begin{longtable}[]{@{}
  >{\raggedright\arraybackslash}p{(\linewidth - 2\tabcolsep) * \real{0.5000}}
  >{\raggedright\arraybackslash}p{(\linewidth - 2\tabcolsep) * \real{0.5000}}@{}}
\toprule\noalign{}
\begin{minipage}[b]{\linewidth}\raggedright
記号
\end{minipage} & \begin{minipage}[b]{\linewidth}\raggedright
経済学的な意味
\end{minipage} \\
\midrule\noalign{}
\endhead
\bottomrule\noalign{}
\endlastfoot
\(\alpha=0\) & 線形生産技術に固定する生産関数の資本弾力性 \\
\(\nu=1\) & 労働不効用水準の正規化定数 \\
\(\beta\in(0,1)\) & 家計の主観的割引因子 \\
\(\gamma>0\) &
異時点間代替弾力性の逆数（\textbf{相対的リスク回避度}） \\
\(\varphi>0\) & フリッシュ労働供給弾力性の逆数 \\
\(\psi_p>1\) & 中間財バラエティ間の代替弾力性 \\
\(\eta_p>0\) & Rotemberg 型価格調整費用パラメータ \\
\(\psi_w>1\) & 差別化された労働タイプ間の代替弾力性 \\
\(\eta_w>0\) & Rotemberg 型賃金調整費用パラメータ \\
\(\tau_p\) & 中間財企業に対する売上比例補助金率 \\
\(\tau_w\) & 労働供給に対する賃金比例補助金率 \\
\(\phi>1\) &
テイラー・ルールにおけるインフレ反応係数（\textbf{テイラー原則}） \\
\(\rho_a,\rho_m\) &
技術ショックおよび金融政策ショックの持続性パラメータ \\
\end{longtable}
}

以下では、フィリップス曲線の傾きを簡潔に表すため、 \[
\kappa_p\equiv\frac{\psi_p}{\eta_p},
\qquad
\kappa_w\equiv\frac{\psi_w}{\eta_w}
\] と定義します。\(\kappa_p\)
は価格マークアップの乖離に対する価格インフレ率の感応度を表し、\(\kappa_w\)
は賃金マークアップの乖離に対する賃金インフレ率の感応度を表します。調整費用パラメータ
\(\eta_p, \eta_w\)
が大きいほど、各インフレ率はマークアップの不均衡に対して緩慢にしか反応しなくなります。

\Needspace{0.18\textheight}

家計の期間効用関数は \[
U(C_t,N_t)
=
\begin{cases}
\dfrac{C_t^{1-\gamma}}{1-\gamma}
-\dfrac{N_t^{1+\varphi}}{1+\varphi},
& \gamma\neq1,\\[6pt]
\log C_t-\dfrac{N_t^{1+\varphi}}{1+\varphi},
& \gamma=1
\end{cases}
\] とします。この効用関数と \(\nu=1\)
のもとでは、消費財単位で測った労働の限界不効用（\textbf{限界代替率：Marginal
Rate of Substitution}）は \[
MRS(C_t,N_t)=C_t^\gamma N_t^\varphi
\]
として与えられます。労働市場の賃金設定では、この限界代替率が労働供給の真の限界費用となります。

\section{モデルの部品}\label{ux30e2ux30c7ux30ebux306eux90e8ux54c1}

\subsection{価格硬直性}\label{ux4fa1ux683cux786cux76f4ux6027}

財市場の価格設定ブロックは第7回と基本的に同一です。中間財企業 \(j\)
は自社製品価格 \(P_t(j)\) を選択し、最終財で測った Rotemberg
型価格調整費用 \(f_P(\Pi_t^p)Y_t\) を支払います。ここで \[
f_P(\Pi_t^p)=\frac{\eta_p}{2}(\Pi_t^p-1)^2,
\qquad
\eta_p>0.
\]
定常状態における独占マークアップの歪みを取り除くため、政府は売上比例是正補助金
\(\tau_p\) を交付します。 \[
1+\tau_p=\frac{1}{1-1/\psi_p}.
\] このとき、ゼロインフレ定常状態において \(MU^p=1\) が保証されます。

対称均衡における非線形価格設定条件は以下の通りです。 \[
\begin{aligned}
\eta_p\Pi_t^p(\Pi_t^p-1)
&=
\psi_p
\left\{
\frac{1}{MU_t^p}
-
(1+\tau_p)\left(1-\frac{1}{\psi_p}\right)
\right\}
\\
&\quad+
\eta_p\mathbb{E}_t
\left[
Q_{t,t+1}
\frac{Y_{t+1}}{Y_t}
\Pi_{t+1}^p(\Pi_{t+1}^p-1)
\right].
\end{aligned}
\] 第7回では実質限界費用 \(MC_t\)
を用いて記述しましたが、本回では後述する賃金ブロックとの対称性を保つため、\textbf{「価格マークアップ（Price
Markup：\(MU_t^p\)）」}を以下のように定義します。 \[
MU_t^p\equiv\frac{\exp(a_t)}{W_t}.
\] 線形生産技術（\(\alpha=0\)）のもとでは労働の限界生産力が
\(\exp(a_t)\) であるため、実質限界費用は
\(MC_t = \frac{W_t}{\exp(a_t)}\)
です。したがって、価格マークアップは実質限界費用の逆数となります。 \[
MU_t^p=\frac{1}{MC_t}.
\] ゼロインフレ定常状態（\(MC=MU^p=1\)）の周りで対数線形化すると、 \[
\mu_t^p=-mc_t
\] が成り立ちます。第7回で導出した価格フィリップス曲線
\(\pi_t^p = \beta\mathbb{E}_t\pi_{t+1}^p + \kappa_p mc_t\)
は、マークアップ表示を用いれば \[
\pi_t^p
=
\beta\mathbb{E}_t\pi_{t+1}^p
-
\kappa_p\mu_t^p
\]
と表現されます。実質限界費用が高まること（\(mc_t>0\)）は、企業の価格マークアップが縮小すること（\(\mu_t^p<0\)）と表裏一体です。

\subsection{労働パッカー}\label{ux52b4ux50cdux30d1ux30c3ux30abux30fc}

賃金硬直性を整合的にモデル化するため、労働サービスは差別化されていると仮定します。連続体
\(h\in[0,1]\) で表される各労働タイプは、それぞれ自らの名目賃金
\(W_t^N(h)\) を独占的に設定します。\textbf{「労働パッカー（Labor
Packer：労働集計業者）」}と呼ばれる完全競争的な事業者は、これらの差別化された個別労働を集計し、同質的な集計労働投入量
\(N_t\) として中間財企業に販売します。

これは、第6回で学んだ財市場の Dixit--Stiglitz
集計を、労働市場へとそのまま鏡写しに適用したものです。労働の集計生産関数は以下のように与えられます。
\[
N_t
=
\left[
\int_0^1
N_t(h)^{\frac{\psi_w-1}{\psi_w}}dh
\right]^{\frac{\psi_w}{\psi_w-1}},
\qquad
\psi_w>1.
\] 労働パッカーは、各個別名目賃金 \(\{W_t^N(h)\}_{h\in[0,1]}\)
と集計名目賃金 \(W_t^N\) を所与として、以下の名目利潤を最大化します。 \[
\max_{\{N_t(h)\}} \left[ W_t^N N_t-\int_0^1 W_t^N(h)N_t(h)dh \right].
\] 一階条件より、労働タイプ \(h\)
に対する需要関数は以下のように導かれます。 \[
N_t(h)
=
N_t
\left(
\frac{W_t^N(h)}{W_t^N}
\right)^{-\psi_w}.
\]
相対的に高い名目賃金を設定した労働タイプに対しては、労働パッカーからの需要が減少します。また、ゼロ利潤条件から得られる集計名目賃金指数は以下の通りです。
\[
W_t^N
=
\left[
\int_0^1
W_t^N(h)^{1-\psi_w}dh
\right]^{\frac{1}{1-\psi_w}}.
\] 全労働タイプが同一の行動をとる対称均衡では、\(W_t^N(h)=W_t^N\) かつ
\(N_t(h)=N_t\)
となります。しかし、各労働タイプは自らが賃金を限界的に引き上げた際に需要が弾力性
\(\psi_w\) で減少することを正確に認識して賃金を決定します。

\subsection{賃金設定と賃金調整費用}\label{ux8cc3ux91d1ux8a2dux5b9aux3068ux8cc3ux91d1ux8abfux6574ux8cbbux7528}

代表的家計は、すべての労働タイプ \(h\in[0,1]\)
の構成員を内包する大家族としてモデル化されます。家計内部には完全な所得保険が存在するため、どの労働タイプであってもすべての構成員は同一の消費水準
\(C_t\) を享受します。一方で、各労働タイプ \(h\) は自らの名目賃金
\(W_t^N(h)\) を自律的に設定します。

賃金設定の動学問題を解くため、労働タイプ \(h\) の労働不効用を
\(\frac{N_t(h)^{1+\varphi}}{1+\varphi}\) と明示します。労働タイプ \(h\)
は、労働パッカーからの需要関数を所与として名目賃金を選択します。政府が賃金比例補助金
\(\tau_w\) を交付するとき、家計が受け取る実効実質賃金は \[
(1+\tau_w)\frac{W_t^N(h)}{P_t}
\]
となります（\textbf{補助金財源の一括税は集計資源配分を歪めないものとします}）。

名目賃金を改定する際、家計は最終財で測った実物的な\textbf{Rotemberg
型賃金調整費用} \[
f_W\left(\frac{W_t^N(h)}{W_{t-1}^N(h)}\right)Y_t,
\qquad
f_W(\Pi_t^w)=\frac{\eta_w}{2}(\Pi_t^w-1)^2,
\quad
\eta_w>0
\] を負担します（\textbf{Rotemberg 1982
の価格調整費用の定式化を賃金設定へと拡張したものです。Calvo
型賃金設定との精緻な対応については Born and Pfeifer 2020, 2025
を参照してください}）。

ここで、\textbf{「賃金マークアップ（Wage
Markup：\(MU_t^w\)）」}を以下のように定義します。 \[
MU_t^w
\equiv
\frac{W_t}{MRS(C_t,N_t)}.
\] 分母の \(MRS(C_t, N_t) = C_t^\gamma N_t^\varphi\)
は労働供給の限界不効用（限界代替率）です。賃金が完全に柔軟であっても、差別化された労働の独占力により、一般にはマークアップ
\(MU_t^w = \frac{1}{(1+\tau_w)(1-1/\psi_w)}\)
が生じます。定常状態の歪みを除く是正補助金率を \[
1+\tau_w=\frac{1}{1-1/\psi_w}
\] と設定することで、ゼロインフレ定常状態において \(MU^w=1\)
が達成されます。賃金が硬直的である場合、現実の賃金マークアップはこの効率的基準から毎期変動します。

労働タイプ \(h\) の相対賃金を \[
s_t^w(h)\equiv \frac{W_t^N(h)}{W_t^N}
\] と定義すると、個別賃金インフレ率は
\(\Pi_t^w(h) = \frac{s_t^w(h)}{s_{t-1}^w(h)}\Pi_t^w\)
と分解されます。前期の相対賃金 \(s_{t-1}^w(h)\)
は今期の賃金調整費用に影響を及ぼすため、動学的な\textbf{状態変数}となります。

労働タイプ \(h\) の価値関数を \(V_t^w(s_{t-1}^w(h))\) とおくと、Bellman
方程式は以下のように定式化されます。 \[
\begin{aligned}
V_t^w(s_{t-1}^w(h))
=
\max_{s_t^w(h)}
\Bigg\{
&
(1+\tau_w)W_t s_t^w(h)N_t(h)
-
\int_0^{N_t(h)}
MRS(C_t,n)\,dn
\\
&
-f_W(\Pi_t^w(h))Y_t
+
\mathbb{E}_t
\left[
Q_{t,t+1}V_{t+1}^w(s_t^w(h))
\right]
\Bigg\}.
\end{aligned}
\] 当期の選択変数 \(s_t^w(h)\)
に関する一階条件および包絡線定理を適用し、対称均衡（\(s_t^w(h)=1, N_t(h)=N_t, MU_t^w(h)=MU_t^w, \Pi_t^w(h)=\Pi_t^w\)）のもとで整理すると、対称均衡における\textbf{非線形賃金設定条件}が得られます。
\[
\begin{aligned}
\eta_w\Pi_t^w(\Pi_t^w-1)
={} &
\psi_w
\left\{
\frac{1}{MU_t^w}
-
(1+\tau_w)\left(1-\frac{1}{\psi_w}\right)
\right\}
\frac{W_tN_t}{Y_t}
\\
&+
\eta_w\mathbb{E}_t
\left[
Q_{t,t+1}
\frac{Y_{t+1}}{Y_t}
\Pi_{t+1}^w(\Pi_{t+1}^w-1)
\right].
\end{aligned}
\]
右辺第1項は現在の賃金マークアップの不均衡を表し、右辺第2項は将来の賃金改定費用を表しています。

\subsection{家計・企業・金融政策}\label{ux5bb6ux8a08ux4f01ux696dux91d1ux878dux653fux7b56}

家計の名目債券保有に関するオイラー方程式は \[
1
=
\mathbb{E}_t
\left[
Q_{t,t+1}
\frac{R_t^N}{\Pi_{t+1}^p}
\right],
\qquad
Q_{t,t+1}
=
\beta
\left(\frac{C_{t+1}}{C_t}\right)^{-\gamma}
\] です。家計の異時点間消費を動かす実質収益率は、価格インフレ率
\(\Pi_{t+1}^p\) で実質化した実質利子率です。

企業は労働パッカーから同質労働 \(N_t\) を雇用して線形生産を行います。 \[
Y_t=\exp(a_t)N_t.
\]
価格調整費用と賃金調整費用が実物資源を消費するため、経済全体の総資源制約式は
\[
Y_t
=
C_t
+
\frac{\eta_p}{2}(\Pi_t^p-1)^2Y_t
+
\frac{\eta_w}{2}(\Pi_t^w-1)^2Y_t
\]
となります。ただし、ゼロインフレ定常状態の周りにおける一次近似では、両方の調整費用は二次の微小量となるため完全に消失します。

中央銀行は、価格インフレ率に反応するテイラー・ルールに従って名目利子率を設定します。
\[
R_t^N
=
\frac{1}{\beta}\exp(-m_t)(\Pi_t^p)^\phi,
\qquad
\phi>1.
\]

\section{非線形体系と定常状態}\label{ux975eux7ddaux5f62ux4f53ux7cfbux3068ux5b9aux5e38ux72b6ux614b}

是正補助金 \(\tau_p^*, \tau_w^*\)
を適用した非線形均衡体系は、以下の連立方程式群にまとめられます。 \[
\begin{aligned}
1
&=
\mathbb{E}_t
\left[
\beta
\left(\frac{C_{t+1}}{C_t}\right)^{-\gamma}
\frac{R_t^N}{\Pi_{t+1}^p}
\right],
\\
MU_t^w
&=
\frac{W_t}{C_t^\gamma N_t^\varphi},
\\
MU_t^p
&=
\frac{\exp(a_t)}{W_t},
\\
\frac{W_t}{W_{t-1}}
&=
\frac{\Pi_t^w}{\Pi_t^p},
\\
Y_t
&=
\exp(a_t)N_t,
\\
Y_t
&=
C_t
+
\frac{\eta_p}{2}(\Pi_t^p-1)^2Y_t
+
\frac{\eta_w}{2}(\Pi_t^w-1)^2Y_t,
\\
R_t^N
&=
\frac{1}{\beta}\exp(-m_t)(\Pi_t^p)^\phi.
\end{aligned}
\]
これに前述の非線形価格設定条件、非線形賃金設定条件、および外生ショック過程を加えることで、モデル全体が完全に閉じます。

\subsection{定常状態}\label{ux5b9aux5e38ux72b6ux614b}

ショックが存在しないゼロインフレ定常状態（\(a=m=0, \Pi^p=\Pi^w=1\)）を導出します。価格調整費用および賃金調整費用がともにゼロとなるため、資源制約式より
\(Y=C\) です。是正補助金により \[
MU^p=MU^w=1
\] が達成されます。

価格マークアップの定義 \(MU^p = \exp(a)/W = 1/W = 1\)
より、定常状態の実質賃金は \[
W=1
\] と定まります。これを賃金マークアップの定義
\(MU^w = W/(C^\gamma N^\varphi) = 1/(N^{\gamma+\varphi}) = 1\)
に代入すると、定常状態の労働投入量は \(N=1\) と求まります。したがって、
\[
N=Y=C=W=1,
\qquad
MU^p=MU^w=1,
\qquad
\Pi^p=\Pi^w=1
\]
という極めて明快な定常状態が得られます。定常状態の粗名目利子率および粗実質利子率は、第7回と同様に
\(R=R^N=1/\beta\) となります。

\section{対数線形化と方程式リスト}\label{ux5bfeux6570ux7ddaux5f62ux5316ux3068ux65b9ux7a0bux5f0fux30eaux30b9ux30c8}

ゼロインフレ定常状態の周りで各均衡方程式を一次近似（\textbf{対数線形化}）します。一次近似のもとでは調整費用の項が落ちるため、財市場の均衡は
\[
y_t=c_t
\] となります。

対数線形化されたモデル全体の方程式体系は、以下の12本の方程式リストに整理されます。

{\def\LTcaptype{none} % do not increment counter
\begin{longtable}[]{@{}
  >{\raggedleft\arraybackslash}p{(\linewidth - 6\tabcolsep) * \real{0.0600}}
  >{\raggedright\arraybackslash}p{(\linewidth - 6\tabcolsep) * \real{0.2400}}
  >{\raggedright\arraybackslash}p{(\linewidth - 6\tabcolsep) * \real{0.4200}}
  >{\raggedright\arraybackslash}p{(\linewidth - 6\tabcolsep) * \real{0.2800}}@{}}
\toprule\noalign{}
\begin{minipage}[b]{\linewidth}\raggedleft
No.
\end{minipage} & \begin{minipage}[b]{\linewidth}\raggedright
方程式の名称
\end{minipage} & \begin{minipage}[b]{\linewidth}\raggedright
方程式の数式表現
\end{minipage} & \begin{minipage}[b]{\linewidth}\raggedright
経済学的な役割
\end{minipage} \\
\midrule\noalign{}
\endhead
\bottomrule\noalign{}
\endlastfoot
1 & フィッシャー方程式 & \(r_t=r_t^N-\mathbb{E}_t\pi^p_{t+1}\) &
事前実質利子率の定義 \\
2 & オイラー方程式 & \(r_t=\gamma\mathbb{E}_t\Delta c_{t+1}\) &
家計の消費の異時点間代替条件 \\
3 & 総資源制約 & \(y_t=c_t\) & 調整費用が消失した財市場均衡 \\
4 & 集計生産関数 & \(y_t=a_t+n_t\) &
線形生産技術のもとでの産出と雇用の連動 \\
5 & 価格マークアップ定義 & \(\mu^p_t=y_t-n_t-w_t=a_t-w_t\) &
労働生産性と実質賃金の乖離（\(-mc_t\)） \\
6 & 価格フィリップス曲線 &
\(\pi^p_t=-\kappa_p\mu^p_t+\beta\mathbb{E}_t\pi^p_{t+1}\) &
企業の価格改定動学（\(\kappa_p \equiv \psi_p/\eta_p\)） \\
7 & 賃金マークアップ定義 & \(\mu^w_t=w_t-\gamma c_t-\varphi n_t\) &
実質賃金と限界代替率（労働供給費用）の乖離 \\
8 & 賃金インフレ恒等式 & \(\pi^w_t=w_t-w_{t-1}+\pi^p_t\) &
名目賃金・物価・実質賃金の変化の定義 \\
9 & 賃金フィリップス曲線 &
\(\pi^w_t=-\kappa_w\mu^w_t+\beta\mathbb{E}_t\pi^w_{t+1}\) &
家計の賃金改定動学（\(\kappa_w \equiv \psi_w/\eta_w\)） \\
10 & テイラー・ルール & \(r_t^N=\phi\pi^p_t-m_t\) &
中央銀行の名目金利設定ルール \\
11 & 技術ショック & \(a_t=\rho_a a_{t-1}+e_t^a\) &
全要素生産性の自己回帰過程 \\
12 & 金融政策ショック & \(m_t=\rho_m m_{t-1}+e_t^m\) &
金融緩和ショックの自己回帰過程 \\
\end{longtable}
}

第8回までのモデルと決定的に異なる点は、賃金インフレ恒等式（No.~8）を通じて、\textbf{内生的な所定状態変数（Endogenous
Predetermined State Variable）としての前期実質賃金
\(w_{t-1}\)}が体系に組み込まれることです。当期の実質賃金 \(w_t\)
は名目賃金インフレ率と価格インフレ率の相互作用によって内生的に決定され、それが翌期への所定状態として引き継がれます。このため、本モデルの動学は単なる外生
AR(1)
ショックへの比例解にはとどまらず、実質賃金の緩慢な調整に伴う持続的な波及効果（\textbf{内生的な慣性}）を伴うことになります。

第7回の価格硬直性のみのモデルへ戻るには、柔軟賃金条件 \(\mu_t^w=0\)
を課せば十分です。このとき No.~7 より \(w_t = \gamma c_t + \varphi n_t\)
となり、No.~8 と No.~9
の賃金ブロックは不要となります。また、価格マークアップは実質限界費用の逆数であるため
\(\mu_t^p = -mc_t\) となり、No.~6
の価格フィリップス曲線は第7回でおなじみの
\(\pi_t^p = \beta\mathbb{E}_t\pi_{t+1}^p + \kappa_p mc_t\)
へと完全に一致します。

\section{マークアップで読む二重の名目硬直性}\label{ux30deux30fcux30afux30a2ux30c3ux30d7ux3067ux8aadux3080ux4e8cux91cdux306eux540dux76eeux786cux76f4ux6027}

第7回では、実質限界費用を産出と技術ショックで表すことで需給ギャップ表現を導出しました。本第10回でも同様の手順を踏みますが、賃金硬直性が加わることで、価格マークアップと賃金マークアップの双方が関与する極めて美しい構造が現れます。

資源制約式 \(c_t=y_t\) と生産関数 \(n_t=y_t-a_t\)
を代入すると、価格マークアップと賃金マークアップはそれぞれ \[
\mu^p_t=a_t-w_t,
\] \[
\mu^w_t
=
w_t-\gamma y_t-\varphi(y_t-a_t)
=
w_t-(\gamma+\varphi)y_t+\varphi a_t
\]
と表されます。ここで、この2つのマークアップを足し合わせると、驚くべきことに\textbf{実質賃金
\(w_t\) が完全に相殺されて消失}します。 \[
\mu^p_t+\mu^w_t
=
(1+\varphi)a_t-(\gamma+\varphi)y_t.
\]
この関係式こそが、二重の名目硬直性モデルを解く最大の鍵です。柔軟価格・柔軟賃金のもとでの自然産出量
\(y_t^f\)
は、両マークアップの歪みがともにゼロ（\(\mu_t^p = \mu_t^w = 0\)）となる水準であるため、
\[
0 = (1+\varphi)a_t-(\gamma+\varphi)y_t^f
\implies
y_t^f = \frac{1+\varphi}{\gamma+\varphi}a_t \equiv \zeta_a a_t,
\qquad
\zeta_a \equiv \frac{1+\varphi}{\gamma+\varphi}
\] となります。これは第7回の自然産出量と完全に一致します。

ここで需給ギャップを \(x_t \equiv y_t - y_t^f\)
と定義すると、マークアップの和は需給ギャップに厳密に比例します。 \[
\mu^p_t+\mu^w_t
=
-(\gamma+\varphi)x_t.
\] したがって、価格マークアップは需給ギャップと賃金マークアップを用いて
\[
\mu^p_t
=
-(\gamma+\varphi)x_t-\mu^w_t
\] と書き換えることができます。

これを価格フィリップス曲線
\(\pi_t^p = \beta\mathbb{E}_t\pi_{t+1}^p - \kappa_p \mu_t^p\)
に代入すると、需給ギャップ表示の価格フィリップス曲線が得られます。 \[
\pi^p_t
=
\beta\mathbb{E}_t\pi^p_{t+1}
+
\kappa_p\{(\gamma+\varphi)x_t+\mu^w_t\}.
\]
この式は極めて重要な経済学的直観を教えてくれます。二重の名目硬直性のもとでは、価格インフレ率は集計的な需給ギャップ
\(x_t\) だけでなく、\textbf{労働市場における賃金マークアップの歪み
\(\mu_t^w\)}によっても直接的に突き動かされるのです。

一方、賃金フィリップス曲線は \[
\pi^w_t
=
\beta\mathbb{E}_t\pi^w_{t+1}
-
\kappa_w\mu^w_t
\]
です。賃金マークアップがプラス（\(\mu_t^w>0\)、すなわち実質賃金が労働供給の限界不効用を上回っている状態）であれば、家計は労働供給を拡大させるために名目賃金を引き下げようとするため、賃金インフレ率には下押し圧力がかかります。

自然利子率は第7回と同様に
\(r_t^f = \gamma\mathbb{E}_t(y^f_{t+1}-y^f_t) = \gamma\zeta_a(\rho_a-1)a_t\)
であり、動学的 IS 曲線も
\(x_t = \mathbb{E}_t x_{t+1} - \frac{1}{\gamma}(r_t-r_t^f)\)
となります。

\subsection{整理形}\label{ux6574ux7406ux5f62}

以上の導出により、二重の名目硬直性をもつ RANK
モデルは、以下の\textbf{5本の整理形連立方程式体系}として美しく集約されます。
\[
\begin{aligned}
x_t
&=
\mathbb{E}_t x_{t+1}
-
\frac{1}{\gamma}(r_t-r_t^f),
\\
\pi^p_t
&=
\beta\mathbb{E}_t\pi^p_{t+1}
+
\kappa_p\{(\gamma+\varphi)x_t+\mu^w_t\},
\\
\pi^w_t
&=
\beta\mathbb{E}_t\pi^w_{t+1}
-
\kappa_w\mu^w_t,
\\
\pi^w_t-\pi^p_t
&=
w_t-w_{t-1},
\\
r_t
&=
\phi\pi^p_t
-
\mathbb{E}_t\pi^p_{t+1}
-
m_t.
\end{aligned}
\] この体系では、賃金マークアップ \(\mu_t^w\) および実質賃金
\(w_t = a_t + (\gamma+\varphi)x_t + \mu_t^w\)
が不可欠な内生変数として連動します。したがって、第7回のように単一の需給ギャップとインフレ率の平面だけでモデルを閉じることはもはやできず、財市場と労働市場の相互連関を同時に追跡する必要があります。

\section{第7回との比較}\label{ux7b2c7ux56deux3068ux306eux6bd4ux8f03}

第7回では価格のみが硬直的で賃金は完全に柔軟であったため、実質賃金は常に労働供給条件
\(w_t = \gamma c_t + \varphi n_t\)
に沿って瞬時に調整され、賃金マークアップは常にゼロ（\(\mu_t^w=0\)）でした。

これに対し本第10回では、名目賃金が硬直的であるため、実質賃金はただちには限界代替率に一致しません。名目賃金がゆっくりとしか改定されない経済において、実質賃金が変化するためには、賃金インフレ率と価格インフレ率の間に\textbf{「乖離（\(\pi_t^w \neq \pi_t^p\)）」}が生じなければなりません。

この違いは、技術ショックに対する反応を考えると極めて鮮明になります。正の技術進歩ショック（\(a_t>0\)）が生じたとき、柔軟価格・柔軟賃金の自然水準では実質賃金が上昇します。第7回の柔軟賃金経済では、名目賃金が瞬時にジャンプして実質賃金の上昇を直ちに達成できました。しかし名目賃金も硬直的である場合、実質賃金を自然水準に向けて引き上げるためには、一定期間にわたって
\[
\pi_t^w > \pi_t^p
\]
という状態を維持し、賃金の上昇ペースが物価の上昇ペースを上回るプロセスを経なければなりません。

\subsection{賃金硬直性のみの場合}\label{ux8cc3ux91d1ux786cux76f4ux6027ux306eux307fux306eux5834ux5408}

理解を深めるため、逆の極限ケースとして「価格は完全に柔軟であり、賃金のみが硬直的である場合」を考察します。価格が完全に柔軟であれば、企業は毎期実質限界費用を定常値に保つため、実質限界費用の対数乖離は常にゼロ（\(mc_t=0\)）、すなわち価格マークアップの乖離も常にゼロ（\(\mu_t^p=0\)）となります。

このとき、マークアップの和の式
\(\mu_t^p + \mu_t^w = -(\gamma+\varphi)x_t\) より、 \[
-\mu_t^w = (\gamma+\varphi)x_t
\] が成り立ちます。これを賃金フィリップス曲線に代入すると、 \[
\pi_t^w
=
\beta\mathbb{E}_t\pi_{t+1}^w
+
\kappa_w(\gamma+\varphi)x_t
\] という、賃金インフレ率に関するおなじみの
NKフィリップス曲線が得られます。

さらに、価格マークアップの定義 \(\mu_t^p = a_t - w_t = 0\)
より、実質賃金は常に技術水準に一致します（\(w_t = a_t\)）。これを賃金インフレ恒等式
\(\pi_t^w - \pi_t^p = w_t - w_{t-1}\) に代入すると、 \[
\pi_t^w - \pi_t^p = a_t - a_{t-1}
\] が得られます。

もし技術ショックが存在しない場合（すべての期で \(a_t=0\)）、上式より \[
\pi_t^w = \pi_t^p
\] が常に成立します。このとき、賃金硬直性のみのモデルは \[
\begin{aligned}
x_t
&=
\mathbb{E}_t x_{t+1}
-
\frac{1}{\gamma}(r_t-r_t^f),
\\
\pi_t
&=
\beta\mathbb{E}_t\pi_{t+1}
+
\kappa_w(\gamma+\varphi)x_t
\end{aligned}
\]
となり、第7回の価格硬直性のみのモデルと\textbf{数学的に全く同一の集計動学体系}へと縮約されます（\(\pi_t \equiv \pi_t^w = \pi_t^p\)）。

ただし、集計量（\(x_t, \pi_t\)）の動学が一致するからといって、経済の内部配分まで一致するわけではありません。純配当を定常消費で規格化した一次偏差
\(d_t = y_t - (w_t + n_t) = -w_t\)
を比較すると、価格硬直性モデルでは柔軟賃金により
\(w_t = (\gamma+\varphi)x_t\) と変動するため純配当は
\(d_t = -(\gamma+\varphi)x_t\)
と大きく振れるのに対し、賃金硬直性モデルでは \(w_t=a_t=0\)
であるため純配当は常に \(d_t=0\)
と一定に保たれます。代表的家計モデル（RANK）では家計が労働所得と配当所得の双方を一括してプールするためこの分配の差異は表面化しませんが、家計を2つに分離する
TANK
モデルでは、この所得分配の帰着がマクロ動学に決定的な差異を生み出すことになります。

また、技術ショックが存在する場合（\(\Delta a_t \neq 0\)）は、実質賃金を動かすために
\(\pi_t^w \neq \pi_t^p\)
とならざるを得ず、この同値性は完全に崩壊します。

\section{まとめと第11回への接続}\label{ux307eux3068ux3081ux3068ux7b2c11ux56deux3078ux306eux63a5ux7d9a}

本講義では、線形生産技術（\(\alpha=0\)）の RANK モデルに労働パッカーと
Rotemberg
型賃金調整費用を導入し、二重の名目硬直性の理論体系を構築しました。得られた核心的結論は以下の通りです。

\begin{enumerate}
\def\labelenumi{\arabic{enumi}.}
\tightlist
\item
  労働市場の独占的競争と賃金調整費用を導入することで、モデルには賃金マークアップ定義、賃金インフレ恒等式、および\textbf{賃金フィリップス曲線}が新たに追加される。
\item
  2つのマークアップを足し合わせると実質賃金が相殺され、\textbf{\(\mu_t^p + \mu_t^w = -(\gamma+\varphi)x_t\)}という要の関係式が得られる。需給ギャップは価格と賃金のマークアップの合計によって分担される。
\item
  名目賃金が硬直的であるため、実質賃金の調整は瞬時ではなく、名目賃金インフレ率と価格インフレ率の差（\(\pi_t^w - \pi_t^p = \Delta w_t\)）を通じて内生的に進行する。前期の実質賃金
  \(w_{t-1}\) は体系の内生的な所定状態変数となる。
\item
  価格と賃金の双方が硬直的な場合、技術ショックによって自然実質賃金が変動すると、価格インフレと賃金インフレの双方をゼロに保ちながら需給ギャップをゼロに維持することは不可能となり、\textbf{神々の配剤は一般に崩壊}する。全期の完全安定化が可能かどうかの詳細な確認については、\hyperref[sec-lecture10-extensions]{発展問題の問5}を参照してください。
\end{enumerate}

続く第11回では、家計の異質性（\textbf{Inequality /
Heterogeneity}）がマクロ経済に及ぼす影響を純粋に抽出するため、賃金を再び柔軟な設定へと戻します。そのうえで、金融資産を保有して異時点間の消費平準化を行う\textbf{「貯蓄家計（Ricardian
Savers）」}と、借入制約に直面して毎期の可処分所得をすべて消費に回す\textbf{「手元流動性制約家計（Hand-to-Mouth
Consumers / Rule-of-Thumb
Households）」}の2つの家計タイプを明示的に導入した\textbf{TANK（Two-Agent
New
Keynesian）モデル}を構築します。第10回で確認した「労働所得と配当所得の分配構造の違い」が、TANK
モデルにおいてまさに政策波及の鍵を握ることになります。

\Needspace{0.35\textheight}

\section{参考文献}\label{ux53c2ux8003ux6587ux732e}

予習は本講義ノートの本文を対象とし、復習では必要に応じて以下の文献の当該論点を参照してください。

\textbf{基本参考文献}

\begin{itemize}
\tightlist
\item
  Erceg, Christopher J., Dale W. Henderson, and Andrew T. Levin (2000)
  ``Optimal Monetary Policy with Staggered Wage and Price Contracts.''
  \emph{Journal of Monetary Economics}, 46(2), 281-313.
  \href{https://doi.org/10.1016/S0304-3932(00)00028-3}{doi}.
  価格と賃金の二重の硬直性のもとで生じる政策トレードオフを包括的に定式化した記念碑的古典。
\item
  Rotemberg, Julio J. (1982) ``Sticky Prices in the United States.''
  \emph{Journal of Political Economy}, 90(6), 1187-1211.
  \href{https://doi.org/10.1086/261117}{doi}.
  価格調整費用の原典。本講義では同一の二次費用関数を賃金設定へと拡張して適用。
\item
  Galí, Jordi (2015) \emph{Monetary Policy, Inflation, and the Business
  Cycle} (2nd ed.). Princeton University Press. 第6章を参照。Calvo
  型の価格・賃金硬直性を RANK モデルへ組み込む標準的体系書。
\item
  Woodford, Michael (2003) \emph{Interest and Prices: Foundations of a
  Theory of Monetary Policy}. Princeton University Press.
  価格・賃金硬直性モデルにおける社会厚生損失関数の二階展開を深く理解するための金字塔。
\end{itemize}

\newpage

\textbf{発展文献}

\begin{itemize}
\tightlist
\item
  Christiano, Lawrence J., Martin Eichenbaum, and Charles L. Evans
  (2005) ``Nominal Rigidities and the Dynamic Effects of a Shock to
  Monetary Policy.'' \emph{Journal of Political Economy}, 113(1), 1-45.
  \href{https://doi.org/10.1086/426038}{doi}.
  価格・賃金の双方の硬直性が金融政策ショックの持続的波及に果たす役割を実証した金字塔的論文。
\item
  Smets, Frank, and Raf Wouters (2003) ``An Estimated Dynamic Stochastic
  General Equilibrium Model of the Euro Area.'' \emph{Journal of the
  European Economic Association}, 1(5), 1123-1175.
  \href{https://doi.org/10.1162/154247603770383415}{doi}.
  二重硬直性を組み込んだ中規模 DSGE モデルの標準推計。
\item
  Calvo, Guillermo A. (1983) ``Staggered Prices in a Utility-Maximizing
  Framework.'' \emph{Journal of Monetary Economics}, 12(3), 383-398.
  \href{https://doi.org/10.1016/0304-3932(83)90060-0}{doi}.
\item
  Svensson, Lars E. O. (2010) ``Inflation Targeting.'' \emph{Handbook of
  Monetary Economics}, Vol. 3B, Chapter 22, 1237-1302.
  \href{https://larseosvensson.se/files/papers/HandbookIT.pdf}{著者版}.
  インフレ・ターゲティングと最適政策の体系的サーベイ。
\item
  Born, Benjamin, and Johannes Pfeifer (2020) ``The New Keynesian Wage
  Phillips Curve: Calvo vs.~Rotemberg.'' \emph{Macroeconomic Dynamics},
  24(5), 1017-1041.
  \href{https://doi.org/10.1017/S1365100518000615}{doi}. Calvo
  型賃金設定と Rotemberg
  型賃金設定の局所的対応関係を精緻に検証した論文。
\item
  Born, Benjamin, and Johannes Pfeifer (2025) ``The New Keynesian Wage
  Phillips Curve: Calvo vs.~Rotemberg -- Corrigendum.''
  \emph{Macroeconomic Dynamics}, 29, e85.
  \href{https://doi.org/10.1017/S1365100525000161}{doi}.
  2020年論文の訂正報告。
\end{itemize}

\Needspace{0.35\textheight}

\section{演習問題}\label{ux6f14ux7fd2ux554fux984c}

\subsection{基本問題}\label{ux57faux672cux554fux984c}

\textbf{問1 {[}Core{]} {[}Derivation{]}：労働パッカーの需要関数導出}

労働集計関数 \[
N_t
=
\left[
\int_0^1
N_t(h)^{\frac{\psi_w-1}{\psi_w}}dh
\right]^{\frac{\psi_w}{\psi_w-1}}
\] のもとで、労働パッカーの利潤最大化問題から、個別労働タイプ \(h\)
に対する需要関数 \[
N_t(h)
=
N_t
\left(
\frac{W_t^N(h)}{W_t^N}
\right)^{-\psi_w}
\] が導出されることを示しなさい。

\textbf{問2 {[}Core{]}：賃金フィリップス曲線の符号の直観}

賃金フィリップス曲線 \[
\pi^w_t
=
\beta\mathbb{E}_t\pi^w_{t+1}
-
\kappa_w\mu^w_t
\] において、賃金マークアップの対数乖離 \(\mu^w_t\)
の係数が負（\(-\kappa_w\)）になる経済学的理由を、労働供給の限界費用と家計の賃金改定インセンティブの観点から説明しなさい。

\textbf{問3 {[}Core{]} {[}Derivation{]}：マークアップの和と需給ギャップ}

財市場の資源制約 \(c_t=y_t\)、生産関数
\(n_t=y_t-a_t\)、およびマークアップの定義式 \[
\mu^p_t=a_t-w_t,
\qquad
\mu^w_t=w_t-\gamma c_t-\varphi n_t
\] を出発点として、実質賃金が相殺されて \[
\mu^p_t+\mu^w_t=-(\gamma+\varphi)x_t
\] が成り立つことを示しなさい。ただし、需給ギャップは
\(x_t \equiv y_t - \frac{1+\varphi}{\gamma+\varphi}a_t\)
と定義されます。

\Needspace{8\baselineskip}

\textbf{問4 {[}Core{]}：名目賃金・物価・実質賃金の動学関係}

実質賃金の時間的変化 \(\Delta w_t \equiv w_t - w_{t-1}\)
を、名目賃金インフレ率 \(\pi_t^w\) と価格インフレ率 \(\pi_t^p\)
を用いて表しなさい。また、経済において実質賃金が上昇するための必要十分条件を述べなさい。

\subsection{発展問題}\label{sec-lecture10-extensions}

問5は本編の動学体系、問6は静学ケースの補論、問7は社会厚生目的関数の補論を用います。

\textbf{問5 {[}Research extension{]}
{[}Derivation{]}：初期実質賃金と完全安定化の不可能性}

価格と賃金がともに硬直的（\(0<\kappa_p,\kappa_w<\infty\)）であり、過去から引き継がれた初期実質賃金
\(w_{-1}\) が所与であるとします。

\begin{enumerate}
\def\labelenumi{\arabic{enumi}.}
\tightlist
\item
  柔軟価格・柔軟賃金の自然配分では自然実質賃金が \(w_t^f=a_t\)
  となることを示し、実質賃金ギャップ \(\widetilde w_t\equiv w_t-a_t\)
  について \[
  \widetilde w_t-\widetilde w_{t-1}
  =\pi_t^w-\pi_t^p-(a_t-a_{t-1})
  \] が成り立つことを導きなさい。
\item
  中央銀行が、時点 \(t=0\) 以降のすべての期において
  \(\pi_t^p=\pi_t^w=0\)
  を達成する完全安定化計画を検討しているとします。両フィリップス曲線から
  \(x_t=0\) かつ \(w_t=a_t\)
  が毎期必要であることを示し、賃金インフレ恒等式と初期条件を連立させて、技術ショックの経路に課される条件を求めなさい。\(a_{-1}=w_{-1}=0\)
  から出発して \(a_0=0.01\)
  となる正の技術ショックが発生した場合、この完全安定化計画は実行可能か論じなさい。
\item
  単に「当期 \(t=0\) だけ \(\pi_0^p=\pi_0^w=0\)
  とする」ことからは、なぜ上記2と同じ結論が出せないのかを説明しなさい。
\end{enumerate}

\Needspace{8\baselineskip}

\textbf{問6 {[}Research extension{]} {[}Derivation{]}
{[}Computation{]}：静学ケースにおける2つのインフレ率と柔軟化極限}

静学補論と同様に、将来期待を
\(\mathbb{E}_t\pi_{t+1}^p=\mathbb{E}_t\pi_{t+1}^w=0\) とおきます。

\begin{enumerate}
\def\labelenumi{\arabic{enumi}.}
\tightlist
\item
  マークアップの和と両フィリップス曲線から \[
  \frac{\pi_t^p}{\kappa_p}+\frac{\pi_t^w}{\kappa_w}
  =(\gamma+\varphi)x_t
  \] を示し、恒等式 \(\pi_t^w-\pi_t^p=\Delta w_t\)
  と連立させて、\(\pi_t^p\) と \(\pi_t^w\)
  の明示的な決定式を求めなさい。
\item
  パラメータを \(\gamma=\varphi=1, \kappa_p=0.10, \kappa_w=0.05\)
  とし、ある期に需給ギャップが \(x_t=0.01\)、実質賃金変化が
  \(\Delta w_t=0.002\)
  であったとします。両インフレ率を計算し、100倍したパーセント表示で示しなさい。需給ギャップがプラス（好況）であるにもかかわらず価格インフレ率がマイナスになり得るメカニズムを、実質賃金調整の観点から説明しなさい。
\item
  実質賃金が変化しない場合（\(\Delta w_t=0\)）の合成傾き
  \(\bar\kappa \equiv (1/\kappa_p + 1/\kappa_w)^{-1}\)
  を求め、賃金調整費用がゼロに近づく極限（\(\eta_w\to0\)）および価格調整費用がゼロに近づく極限（\(\eta_p\to0\)）において
  \(\bar\kappa\) がどのような値に収束するかを調べなさい。
\item
  この静学的な関係式だけでは、技術ショックや金融政策ショックに対する動学的なインパルス応答を計算したことにならない理由を述べなさい。
\end{enumerate}

\Needspace{20\baselineskip}

\textbf{問7 {[}Research extension{]}
{[}Derivation{]}：実質賃金調整に伴う名目調整費用の最小化}

目的関数補論で導出した社会厚生ウェイト
\(\omega_p \equiv \frac{\eta_p}{\gamma+\varphi}>0, \omega_w \equiv \frac{\eta_w}{\gamma+\varphi}>0\)
を用います。必要な当期の実質賃金変化が外生的に \(\Delta w_t=\delta\)
として与えられているとき、名目調整費用に起因する当期の損失 \[
J_t=\frac12\left[\omega_p(\pi_t^p)^2+\omega_w(\pi_t^w)^2\right]
\] を考えます。

\begin{enumerate}
\def\labelenumi{\arabic{enumi}.}
\tightlist
\item
  恒等式 \(\pi_t^w-\pi_t^p=\delta\) を制約条件として \(J_t\)
  を最小化し、最適な価格インフレ率 \(\pi_t^p\)、賃金インフレ率
  \(\pi_t^w\)、および損失の最小値を求めなさい。\(\delta>0\)
  のときの符号と、賃金の調整費用ウェイト \(\omega_w\)
  が相対的に極めて大きい場合に調整負担がどのように分担されるかを論じなさい。
\item
  \(\gamma=\varphi=1, \eta_p=60, \eta_w=120, \delta=0.01\)
  の場合の数値を計算しなさい。ここで得られた値が、動学的一般均衡モデル全体における最適政策の厚生損失そのものではない理由を説明しなさい。
\item
  \(\delta\ne0\) のとき、上記1で求めた最適なインフレ率が、問6の
  \(x_t=0, \Delta w_t=\delta\)
  のもとでの静学均衡解と完全に一致するための条件を求め、価格と賃金の代替弾力性
  \(\psi_p, \psi_w\) の関係として表現しなさい。
\end{enumerate}

\clearpage

\section{補論（発展）：期待項を落とした静学ケース}\label{ux88dcux8ad6ux767aux5c55ux671fux5f85ux9805ux3092ux843dux3068ux3057ux305fux9759ux5b66ux30b1ux30fcux30b9}

この補論は、動学モデルにおける前向きな期待項を一旦捨象することで、価格マークアップと賃金マークアップが需給ギャップをどのように分担しているかを直観的に確認するための発展項目です。

本来のニューケインジアン・モデルは動学体系であり将来期待と同時に解く必要がありますが、代数構造の美しさを際立たせるため、将来インフレ期待がゼロである仮想的な静学ケース
\[
\mathbb{E}_t\pi_{t+1}^p=\mathbb{E}_t\pi_{t+1}^w=0
\]
を考えます。このとき、価格および賃金のフィリップス曲線は以下のように極めて簡潔になります。
\[
\pi_t^p
=
-\kappa_p\mu_t^p,
\qquad
\pi_t^w
=
-\kappa_w\mu_t^w.
\] ここにマークアップの和の基本式
\(\mu_t^p+\mu_t^w = -(\gamma+\varphi)x_t\) を代入すると、 \[
\frac{\pi_t^p}{\kappa_p}
+
\frac{\pi_t^w}{\kappa_w}
=
(\gamma+\varphi)x_t
\] という関係式が得られます。この式は、「需給ギャップ \(x_t\)
が生み出すインフレ圧力を、価格インフレ率と賃金インフレ率がそれぞれの価格改定の柔軟性（感応度
\(\kappa_p, \kappa_w\)）に応じて分担して吸収している」という事実を端的に物語っています。

さらに、実質賃金の変化に関する恒等式
\(\Delta w_t \equiv w_t - w_{t-1} = \pi_t^w - \pi_t^p\) を連立させて
\(\pi_t^p\) と \(\pi_t^w\) について解くと、 \[
\pi_t^p
=
\frac{\kappa_p\kappa_w}{\kappa_p+\kappa_w}
(\gamma+\varphi)x_t
-
\frac{\kappa_p}{\kappa_p+\kappa_w}\Delta w_t,
\] \[
\pi_t^w
=
\frac{\kappa_p\kappa_w}{\kappa_p+\kappa_w}
(\gamma+\varphi)x_t
+
\frac{\kappa_w}{\kappa_p+\kappa_w}\Delta w_t
\]
が得られます。この2式から二重硬直性の本質が明瞭に理解できます。実質賃金を引き上げる必要がある局面（\(\Delta w_t>0\)）では、賃金インフレ率は価格インフレ率よりも必ず高くなければなりません。

特に、実質賃金が変化しない局面（\(\Delta w_t=0\)）を考えると、両インフレ率は一致し、
\[
\pi_t^p=\pi_t^w
=
\bar{\kappa}(\gamma+\varphi)x_t,
\qquad
\bar{\kappa}
\equiv
\left(
\frac{1}{\kappa_p}
+
\frac{1}{\kappa_w}
\right)^{-1}
\] となります。ここで \(\bar{\kappa}\) は、\(\kappa_p\) と \(\kappa_w\)
の\textbf{調和平均型の合成傾き係数}です。例えば \(\kappa_p=\kappa_w\)
であれば、\(\bar{\kappa} = \kappa_p/2\)
となり、フィリップス曲線の傾きは価格硬直性のみの場合の半分へと平坦化します。

図~\ref{fig-lecture10-static-slopes}
は、この静学ケースにおいて、価格と賃金の二重の硬直性がフィリップス曲線の傾きをいかに緩やかにするかを幾何学的に示しています。

\begin{figure}

\centering{

\includegraphics[width=0.9\linewidth,height=\textheight,keepaspectratio]{lecture10_files/mediabag/figures/lecture10/lecture10_static_case_slopes.pdf}

}

\caption{\label{fig-lecture10-static-slopes}期待項を落とした静学ケースの傾き比較：\(\mathbb{E}_t\pi_{t+1}^p=\mathbb{E}_t\pi_{t+1}^w=0\)
かつ \(\Delta w_t=0\)
のもとで、価格硬直性のみ、賃金硬直性のみ、および二重硬直性のフィリップス曲線の傾きを比較した図。二重硬直性の合成傾きは調和平均
\(\bar\kappa(\gamma+\varphi)\) となり、最も平坦化する。}

\end{figure}%

\section{補論（発展）：目的関数の導出}\label{ux88dcux8ad6ux767aux5c55ux76eeux7684ux95a2ux6570ux306eux5c0eux51fa}

本補論では、第9回と同様のアプローチを用い、代表的家計の効用関数と総資源制約式を定常状態の周りで二階展開することにより、中央銀行が最小化すべき社会厚生損失関数を厳密に導出します。

本編と同様に
\(\alpha=0\)、政府支出なし、\(a=0\)、\(\nu=1\)、および定常状態水準
\(C=N=Y=W=1, MU^p=MU^w=1\) を基準とします。

\subsection{資源制約の二階近似}\label{ux8cc7ux6e90ux5236ux7d04ux306eux4e8cux968eux8fd1ux4f3c}

価格調整費用および賃金調整費用を含む総資源制約式は \[
Y_t
=
C_t
+
\frac{\eta_p}{2}(\Pi_t^p-1)^2Y_t
+
\frac{\eta_w}{2}(\Pi_t^w-1)^2Y_t
\] です。両辺を \(Y_t\)
で割り、消費財として利用可能な割合について解くと、 \[
C_t
=
Y_t
\left[
1
-
\frac{\eta_p}{2}(\Pi_t^p-1)^2
-
\frac{\eta_w}{2}(\Pi_t^w-1)^2
\right]
\] となります。粗インフレ率の対数展開
\(\Pi_t-1 = \pi_t + \frac{1}{2}\pi_t^2 + O(3)\)
より、調整費用の二階近似は \[
\Phi_t
\equiv
\frac{\eta_p}{2}(\pi_t^p)^2
+
\frac{\eta_w}{2}(\pi_t^w)^2
\]
と定義されます。資源制約の両辺に対数をとってテイラー展開すると、二階の精度で
\[
c_t
=
y_t
-
\frac{\eta_p}{2}(\pi_t^p)^2
-
\frac{\eta_w}{2}(\pi_t^w)^2
+O(3)
\]
が得られます。価格調整費用と賃金調整費用は、実物資源の直接的な散逸を通じて消費水準を押し下げます。

\Needspace{0.28\textheight}

\subsection{効用関数の二階近似}\label{ux52b9ux7528ux95a2ux6570ux306eux4e8cux968eux8fd1ux4f3c}

期間効用関数
\(U(C_t, N_t) = \frac{C_t^{1-\gamma}-1}{1-\gamma} - \frac{N_t^{1+\varphi}}{1+\varphi}\)
を定常状態の周りで対数二階近似します。\(Z_t = \exp(z_t) = 1 + z_t + \frac{1}{2}z_t^2 + O(3)\)
を適用すると、 \[
U(C_t,N_t)
=
\text{const.}
+
c_t-n_t
+
\frac{1}{2}
\left[
(1-\gamma)c_t^2
-
(1+\varphi)n_t^2
\right]
+O(3)
\] が得られます。線形項の \(c_t\) に資源制約式 \(c_t = y_t - \Phi_t\)
を代入し、生産関数 \(n_t = y_t - a_t\) を用いて整理すると、 \[
U(C_t,N_t)
=
\text{const.}
+a_t
-\Phi_t
+
\frac{1}{2}
\left[
-(\gamma+\varphi)y_t^2
+2(1+\varphi)a_t y_t
-(1+\varphi)a_t^2
\right]
+O(3)
\]
となります。政策運営に依存しない外生ショックのみの項を落とし、産出に関する二次形式を平方完成すると、
\[
U(C_t,N_t)
=
\text{policy-independent terms}
-
\frac{\gamma+\varphi}{2}
\left(
y_t-\frac{1+\varphi}{\gamma+\varphi}a_t
\right)^2
-
\Phi_t
+O(3)
\] が得られます。

自然産出量 \(y_t^f = \frac{1+\varphi}{\gamma+\varphi}a_t\)
および需給ギャップ \(x_t \equiv y_t - y_t^f\) を代入すると、期間効用は
\[
\begin{aligned}
U(C_t,N_t)
&=
\text{policy-independent terms}
\\
&\quad-
\frac{1}{2}
\left[
(\gamma+\varphi)x_t^2
+
\eta_p(\pi_t^p)^2
+
\eta_w(\pi_t^w)^2
\right]
+O(3)
\end{aligned}
\] として美しく導出されます。

\subsection{目的関数}\label{ux76eeux7684ux95a2ux6570}

したがって、中央銀行が最小化すべき期ごとの社会厚生損失関数は \[
\ell_t
=
\frac{1}{2}
\left[
(\gamma+\varphi)x_t^2
+
\eta_p(\pi_t^p)^2
+
\eta_w(\pi_t^w)^2
\right]
\] となります。全体を正の定数 \(\gamma+\varphi\) で割って正規化すると、
\[
\ell_t
=
\frac{1}{2}
\left[
x_t^2
+
\omega_p(\pi_t^p)^2
+
\omega_w(\pi_t^w)^2
\right],
\qquad
\omega_p
=
\frac{\eta_p}{\gamma+\varphi},
\qquad
\omega_w
=
\frac{\eta_w}{\gamma+\varphi}
\]
が得られます。第9回の損失関数に比べて、賃金調整費用の存在を反映する\textbf{賃金インフレ率の二乗項
\(\omega_w(\pi_t^w)^2\)}が新たに追加されていることが確認できます。




\end{document}
